\scanpage{68}
% archive sheet number: 67
contour orienté, et
\[
\underset{\substack{\shortparallel\\Y_{\gamma(\sigma)}}}{Y_\sigma}
=\prod_{i\in I_1(\sigma)}Y_i
\left(\prod_{a\in I_2(\sigma)/\sigma_2}Y_a\right)^2
\prod_{a\in\underset{\substack{|\\\text{arêtes}}}{A}(I_3(\sigma))}Y_a
\]
Quant aux monômes \inserted{$Y_\sigma$}, ils correspondent $1$-$1$ aux structures de graphes \deleted{\unsure{à}} stricts \inserted{$\gamma$} sur $I$ correspondants aux $\sigma$, i.e dont les composantes connexes sont soit des sommets isolés, soit des \deleted{simplets} doublets
\figure{\begin{tikzpicture}[baseline=-.5ex]
\fill (0,0) circle (1.3pt) (.8,0) circle (1.3pt);
\draw (0,0)--(.8,0);
\end{tikzpicture}}
, soit des polygônes combinatoires d'ordre \deleted{\unsure{et|d}} $\geq3$. Pour que $Y_\sigma=Y_{\sigma'}$, il f. et s. \inserted{donc} que $\sigma$ et $\sigma'$ définissent les $=$ $I_1$ et $I_2$, et \deleted{que $I_1(\sigma)$} la $=$ structure \deleted{polyg.} de contour sur $I_3=I\setminus I_1\setminus I_2$, donc qu'ils \deleted{\unsure{le prod.|la partie}} définissent $=$ partition de $I_3$ en \deleted{parties connexes} \inserted{orbites et} que sur chaque classe $\sigma$ et $\sigma'$ sont ou bien égales, ou bien inverses. Donc chaque $Y_{\gamma\deleted{\unsure{\sigma}}}$ intervient exactement $2^{\unsure{\nu|v}_*}$ fois. Donc
\[
\begin{aligned}
\Delta={}&\sum_{\gamma\in\underset{\substack{\downarrow\\\text{ens des structures de}\\\text{graphes sur }I\text{ du type}\\\text{envisagé}}}{\Gamma_I}}
(-1)^{\nu_2(\gamma)+\nu_*^+(\gamma)}
2^{\overbrace{\nu_1(\gamma)+\nu_*(\gamma)}}\,
\underbrace{\begin{aligned}
&\prod_{i\in I_1(\unsure{\sigma|\gamma})}Y_i
\left(\prod_{a\in A(I_2(\gamma))}Y_a\right)^2\\
&\qquad\prod_{a\in A(I_3(\gamma))}Y_a
\end{aligned}}_{M_\gamma}
\end{aligned}
\]

